% ---- pg-061 （书页 49） ----

设 $k\times k$ 对称矩阵 $\bm{A}$ 对于一切 $\bm{x}'=[x_1,x_2,\cdots,x_k]$，满足

\phantomsection\label{eq:2-17}
\begin{equation}
0\leqslant \bm{x}'\bm{A}\bm{x} \tag{2-17}
\end{equation}

则称 $\bm{A}$ 是非负定的.若仅当向量 $\bm{x}'=[0,0,\cdots,0]$ 时，\hyperref[eq:2-17]{式(2-17)}中等号才成立，则称 $\bm{A}$ 是正定的.换言之，如果对一切 $\bm{x}\neq \bm{0}$ 的向量，有

\phantomsection\label{eq:2-18}
\begin{equation}
0<\bm{x}'\bm{A}\bm{x} \tag{2-18}
\end{equation}

则称 $\bm{A}$ 是正定的.因为 $\bm{x}'\bm{A}\bm{x}$ 中只含有平方项 $x_i^2$ 及乘积项 $x_ix_j$，故称为二次型.

\mvbold{例 2.11} \quad (正定矩阵和二次型)

证明以下二次型是正定的：

\[
3x_1^2+2x_2^2-2\sqrt{2}x_1x_2
\]

为了说明一般方法，首先写出二次型的矩阵表示形式

\[
[x_1 \quad x_2]
\begin{bmatrix}
3 & -\sqrt{2}\\
-\sqrt{2} & 2
\end{bmatrix}
\begin{bmatrix}
x_1\\
x_2
\end{bmatrix}
=\bm{x}'\bm{A}\bm{x}
\]

由定义 2A.30，$\bm{A}$ 的特征值是方程 $|\bm{A}-\lambda \bm{I}|=0$ 或 $(3-\lambda)(2-\lambda)-2=0$ 的解.求解得，$\lambda_1=4$，$\lambda_2=1$.利用\hyperref[eq:2-16]{式(2-16)}的谱分解，有

\[
\underset{(2\times2)}{\bm{A}}=\underset{(2\times1)}{\lambda_1\bm{e}_1}\underset{(1\times2)}{\bm{e}_1'}+\underset{(2\times1)}{\lambda_2\bm{e}_2}\underset{(1\times2)}{\bm{e}_2'}
\]
\[
=\underset{(2\times1)}{4\bm{e}_1}\underset{(1\times2)}{\bm{e}_1'}+\underset{(2\times1)}{\bm{e}_2}\underset{(1\times2)}{\bm{e}_2'}
\]

其中 $\bm{e}_1$ 和 $\bm{e}_2$ 是分别对应于特征值 $\lambda_1=4$ 和 $\lambda_2=1$ 的标准化的正交的特征向量.因为 $4$ 和 $1$ 是标量，用 $\bm{x}'$ 和 $\bm{x}$ 分别左乘和右乘 $\bm{A}$，其中 $\bm{x}'=[x_1,x_2]$ 是任一非零向量，得

\[
\underset{(1\times2)}{\bm{x}'}\underset{(2\times2)}{\bm{A}}\underset{(2\times1)}{\bm{x}}=\underset{(1\times2)}{4\bm{x}'}\underset{(2\times1)}{\bm{e}_1}\underset{(1\times2)}{\bm{e}_1'}\underset{(2\times1)}{\bm{x}}+\underset{(1\times2)}{\bm{x}'}\underset{(2\times1)}{\bm{e}_2}\underset{(1\times2)}{\bm{e}_2'}\underset{(2\times1)}{\bm{x}}
\]
\[
=4y_1^2+y_2^2\geqslant 0
\]

其中

\[
y_1=\bm{x}'\bm{e}_1=\bm{e}_1'\bm{x} \qquad \text{和} \qquad y_2=\bm{x}'\bm{e}_2=\bm{e}_2'\bm{x}
\]

这表明，$y_1$ 和 $y_2$ 不全为零，因此，$\bm{x}'\bm{A}\bm{x}=4y_1^2+y_2^2>0$，即 $\bm{A}$ 是正定的.

从 $y_1$ 和 $y_2$ 的定义，有

\[
\begin{bmatrix}
y_1\\
y_2
\end{bmatrix}
=
\begin{bmatrix}
\bm{e}_1'\\
\bm{e}_2'
\end{bmatrix}
\begin{bmatrix}
x_1\\
x_2
\end{bmatrix}
\]

或

\[
\underset{(2\times1)}{\bm{y}}=\underset{(2\times2)}{\bm{E}}\underset{(2\times1)}{\bm{x}}
\]

现在 $\bm{E}$ 是正交阵，因而有逆矩阵 $\bm{E}'$.于是，$\bm{x}=\bm{E}'\bm{y}$.$\bm{x}$ 是非零向量，则 $\bm{0}\neq \bm{x}=\bm{E}'\bm{y}$ 意味着 $\bm{y}\neq \bm{0}$.

运用谱分解，很容易证明，当且仅当 $\bm{A}$ 的所有特征值均大于零时，$k\times k$ 对称阵 $\bm{A}$ 是正定矩阵(参考练习 2.17).当且仅当 $\bm{A}$ 的所有特征值大于或等于零时，$\bm{A}$ 是非负定阵.

现在假定向量 $\bm{x}$ 的 $p$ 个元素 $x_1,x_2,\cdots,x_p$ 被视为 $p$ 个随机变量 $X_1,X_2,\cdots,X_p$.正如第 1 章中指出的，可以把这些元素当做 $p$ 维空间中一个点的坐标；在这种情况下，点 $(x_1,x_2,\cdots,x_p)$ 到原点的“距离”应当利用单位标准偏差来表示.这样，我们可以解释观察结果中所固有的不确定性(变化性).具有相同的“不确定性”的点，可以看成与原点的距离相同.

利用第 1 章中介绍的距离公式[参见\hyperref[eq:1-22]{式(1-22)}]，假定对所有的 $[x_1,x_2,\cdots,x_p]\neq[0,0,\cdots,0]$，有(距离)$^2>0$，则到原点的距离满足一般公式
